\section{Recopilación y aseguramiento de calidad de los datos de demostración}

\subsection{Pipeline de recopilación de datos}

Los datos de demostración provienen de humanos, robots simulados o pre-recorded log. Pipelines comunes:
\begin{itemize}
    \item \textbf{Kinesthetic Teaching}: el cuerpo humano controla directamente el brazo robótico, con baja latency.
    \item \textbf{Remote teleoperation}: control en cuadrícula, que permite expertos distribuidos geográficamente.
    \item \textbf{Behavioral cloning from screen recording}: sincronización de video con la actividad de teclado y mouse.
\end{itemize}

\begin{knowledgebox}{embodiment gap}
Usar directamente video humano en lugar de robot demonstration produce un embodiment gap (diferencia de grados de libertad, diferencia de dinámica), pero puede servir como exploration signal, y el gap se puede reducir mediante domain randomization.
\end{knowledgebox}

\subsection{QA de datos y audit}

Las demostraciones de alta calidad requieren un audit pipeline, que incluye:
\begin{itemize}
    \item agreement rate entre expertos (Cohen's κ);
    \item variance coverage de los pares estado-acción;
    \item labeling de las demostraciones erróneas;
    \item Metadata (dificultad de la tarea, lighting del entorno).
\end{itemize}

\begin{table}[H]
\centering
\begin{tabular}{p{0.28\textwidth} p{0.35\textwidth} p{0.33\textwidth}}
\toprule
Dimensión & Métrica & Respuesta de ingeniería \\
\midrule
Consistencia & Pairwise agreement & Los ejemplos inconsistentes se envían a revisión \\
Diversidad & Different skills, contexts & Se usa clustering para identificar gaps de cobertura \\
Seguridad & Adversarial prompt & Se agregan counterfactual data \\
\bottomrule
\end{tabular}
\caption{Aseguramiento de calidad de los datos de demostración}
\end{table}

\subsection{Colaboración entre modalidades y entre roles}

Los proyectos modernos de aprendizaje por imitación suelen involucrar a equipos de etiquetado de datos, expertos en lenguaje y SRE, que mantienen en conjunto el dataset catalog y el manifest. Cada rollout requiere agregar new prompts al dataset y registrar el shift log en el evidence board.

\subsection{Resumen de la sección}

La calidad de los datos de demostración determina directamente el límite inferior del aprendizaje por imitación. Mediante un audit riguroso, metadata y multi-role feedback, es posible convertir problemas como el embodiment gap o el adversarial prompt en riesgos gobernables.

